\documentclass[11pt]{article}
\usepackage{amsmath, amssymb, amsthm}

\title{Category Theory, Index Theory, and Positive Geometry:\\
A Structural Correspondence Without False Equivalences}
\author{Borealis S. Hedling}
\date{}

\begin{document}
\maketitle

\section*{Abstract}
This document provides a mathematically consistent relationship between three
independent domains: category theory (via the Yoneda Lemma), index theory (via
the Atiyah--Singer Index Theorem), and positive geometry (via the amplituhedron).
Unlike informal diagrams that imply direct morphisms between these fields, we
construct a structural correspondence that respects the mathematical boundaries
of each domain. No claim is made of a universal functor or ternary encoding
linking these theories; instead, we identify the correct conceptual interfaces
that allow them to coexist in a unified meta-framework.

\section{Category-Theoretic Structure}
Let $\mathcal{C}$ be a locally small category. The Yoneda Lemma states that
for any objects $X,Y \in \mathcal{C}$,
\[
\mathrm{Nat}(\mathrm{Hom}(-,X), \mathrm{Hom}(-,Y))
\cong \mathrm{Hom}(X,Y).
\]
This establishes an equivalence between natural transformations and morphisms,
but does \emph{not} imply any relationship to geometric scattering amplitudes
or index-theoretic invariants.

\section{Index-Theoretic Structure}
Let $D$ be an elliptic differential operator on a compact manifold $M$ with
vector bundle $E$. The Atiyah--Singer Index Theorem states:
\[
\mathrm{Index}(D)
= \int_M \hat{A}(M)\,\mathrm{ch}(E).
\]
This index is a topological invariant and cannot be encoded in a discrete
ternary alphabet $\{-1,0,+1\}$ without loss of structure. No canonical functor
exists from elliptic operators to discrete computational encodings.

\section{Positive Geometry and the Amplituhedron}
The amplituhedron $\mathcal{A}_{n,k,m}$ is a positive geometry defined by the
image of the positive Grassmannian $G_+(k,n)$ under a linear map into
$\mathbb{R}^{k(m+k)}$. It encodes scattering amplitudes in planar $\mathcal{N}=4$
super Yang--Mills theory. Formally,
\[
\mathcal{A}_{n,k,m} = \Phi(G_+(k,n)),
\]
where $\Phi$ is a positivity-preserving map. This geometry is unrelated to
category-theoretic Hom-functors or index-theoretic integrals.

\section{A Correct Structural Correspondence}
Although these domains do not admit direct morphisms, they share a common
meta-structure: each provides a method for extracting invariant information
from complex systems.

\begin{itemize}
    \item Category theory extracts \emph{structural invariants} via natural
    transformations.
    \item Index theory extracts \emph{topological invariants} via elliptic
    operators.
    \item Positive geometry extracts \emph{scattering invariants} via
    canonical forms.
\end{itemize}

We define a meta-framework:
\[
\mathcal{S} = \{\text{structural},\ \text{topological},\ \text{geometric}\}
\]
and a correspondence map
\[
\Xi : \{\text{Yoneda},\ \text{Index},\ \text{Amplituhedron}\} \to \mathcal{S},
\]
which assigns each theory to its invariant-extraction role. This avoids false
equivalences while preserving conceptual unity.

\section{Conclusion}
There is no universal functor, no ternary encoding, and no direct morphism
connecting the Yoneda Lemma, the Atiyah--Singer Index Theorem, and the
amplituhedron. The corrected framework presented here respects the mathematical
independence of each domain while identifying their shared role in extracting
invariant structure from complex systems.

\end{document}
